\section{Introduction}

Cellular automata are discrete dynamical systems that have been used to model complex systems~\cite{Torres2025_complejo,sole1996_cs_caos,Hoekstra2010_cs_ca}: from very simple local rules they produce behaviors that range from convergence to a homogeneous, easily predictable state to chaotic dynamics~\cite{wolfram1982_ca,Wolfram2002_nks}. To distinguish these behaviors, several classifications have been proposed, such as Wolfram's four classes~\cite{Wolfram2002_nks}, the classification of Li and Packard~\cite{li1990_structure} and others reviewed by Martínez~\cite{genaro2013_class}, or the set of \emph{life-like} rules~\cite{lifewiki_life}. But how similar are the rules within the same class? And how different are they from the rules of another class?

Apart from directly comparing their exact behavior, classifications are the only measure to decide whether two rules are alike, and both give a binary answer. The purpose of this research is to define a method to compare the behavior of any two rules, even if they have different numbers of states or neighbors, and to measure their degree of similarity.

We first identify the exact relations that exist between rules of different domains (Section~\ref{sec:inheritance}) and use them as a reference to validate the proposed methods. We then compare the transition graph through its canonical form (Section~\ref{sec:structural}) and propose spectral methods that represent each rule by a vector (Section~\ref{sec:spectral}): the eigenvalues, the spectral moments and the spectral moments with attractor mass. Finally, we evaluate this last method on rules from different domains, on spaces of different sizes, and by clustering the elementary rules (Section~\ref{sec:experiments}).
